\section{Related Work}
\label{sec:related-work}

\subsection{Executable Software-Engineering Benchmarks}

Executable software-engineering benchmarks have established the importance of
real repositories, tool use, sandboxed execution, and test-based verification.
SWE-bench asks language models to resolve real GitHub issues in existing
repositories~\cite{jimenez2024swebench}, while Terminal-Bench evaluates agents
on realistic command-line tasks~\cite{merrill2026terminalbench}. Their central
unit is typically an issue, a failing test, or a requested modification to the
original workspace. \flb instead evaluates extraction into a new independent
artifact. The source repository is evidence, not the final execution
environment, and passing requires behavior preservation after the source
package is unavailable.

Agent-computer interfaces and runtime design can also change executable-agent
performance~\cite{yang2024sweagent,yao2026harnessbench}. We therefore fix
OpenHands as the \mainprotocol runtime for the model leaderboard. Runtime
comparisons, when available, are reported separately rather than mixed into
model scores.

\subsection{Code Generation, Repository Understanding, and Localization}

Function-level code generation measures whether a model can synthesize a
program from a natural-language or executable specification, as in HumanEval
and related synthesis
suites~\cite{chen2021codex,austin2021mbpp,lai2023ds1000,zhuo2024bigcodebench,jain2024livecodebench}.
Repository-level completion and retrieval study whether a model can use
in-repo context to finish or locate code~\cite{zhang2023repocoder,liu2024repobench}.
Feature-location research asks whether a developer or tool can identify files,
entities, or regions relevant to a concern~\cite{dit2013featurelocation}.
Feature lifting requires both inference and artifact construction, but success
is not reducible to either. Locating the correct implementation is insufficient
if transitive behavior or resources are omitted; generating a behaviorally
similar implementation is insufficient if the public contract or independent
package boundary is not preserved.

\subsection{Program Slicing, Modularization, and Library Extraction}

Feature extraction is related to program slicing~\cite{weiser1984slicing},
multi-dimensional separation of concerns~\cite{tarr1999degrees}, and the
refactoring literature that surveys how developers extract and modularize
behavior~\cite{mens2004refactoring}. Classical techniques typically assume a
program representation, slicing criterion, or human-specified boundary. \flb
evaluates an autonomous agent from a natural-language public contract and a
complete repository, and grades the final independent artifact by observable
behavior. The benchmark does not claim to compute a minimal semantic slice: its
frozen references are feasible comparison points, and \rres is a compactness
proxy rather than a proof of optimality.

\subsection{Information Boundaries in Agent Evaluation}

Agent performance depends on the model, execution runtime, tools, context
policy, and feedback interface~\cite{yang2024sweagent,wang2024openhands,yao2026harnessbench}.
\flb separates these factors in two ways. The main leaderboard fixes OpenHands
as the \mainprotocol runtime and varies the model backend. Information
ablations instead change what task evidence is visible while preserving model,
runtime, evaluator, and budget. This design distinguishes model comparison from
questions such as whether executable Public tests or source-location hints
alter performance. Runtime replications, when completed, are reported
separately rather than mixed into model scores.

